\documentclass{book}

\usepackage[utf8]{inputenc}  % para que funcionen las tildes
\usepackage{amsmath}
\usepackage{amssymb}
\usepackage{amsthm}
\usepackage{letltxmacro}
\LetLtxMacro\amsproof\proof
\LetLtxMacro\amsendproof\endproof
\usepackage{txfonts}
\usepackage{lmodern}
\usepackage[dvipsnames]{xcolor}
\usepackage{thmtools, thm-restate}
\usepackage{datetime} % para la hora de compilación
\usepackage{graphicx}
\graphicspath{{images/}}
\usepackage[spanish,es-noquoting]{babel} % es-noquoting es para que funcione tikz
\usepackage{mathabx} % para \divides
\usepackage{centernot} % para \centernot\inculdes
\usepackage{wrapfig}
\usepackage{multicol}
\usepackage{subcaption}
\usepackage{xfrac}
\usepackage{glossaries}
\usepackage{standalone}
\usepackage{pgfmath}
\usepackage{pgfplots}
\usepackage{pdfpages}
\usepackage{tikz}
\usepackage{mathtools}
\usepackage{hyperref}
\usetikzlibrary{arrows.meta}
\usetikzlibrary{shapes}
\usetikzlibrary{decorations.text}
\usetikzlibrary{positioning}
\usetikzlibrary{external}
\tikzexternalize[prefix=./tikzbuild/]
\tikzexternalize % activate!

\AtBeginDocument{%
  \LetLtxMacro\proof\amsproof
  \LetLtxMacro\endproof\amsendproof
}

\definecolor{coolblack}{rgb}{0.0, 0.18, 0.39}
%\usepackage[a6paper,margin=5mm]{geometry}
\usepackage[a4paper, top=2.5cm,bottom=2.5cm,left=2.5cm,right=2.5cm]{geometry}

%\setlength{\parindent}{0pt}
\usepackage{parskip}

\usepackage{hyperref}

% ESTILOS DE DEFINICIONES Y TEOREMAS
\declaretheoremstyle[
	bodyfont=\normalfont,
	shaded={
		margin=8pt,
		bgcolor=White,
		rulecolor=Black,
		rulewidth=1pt
}]{mythm}

\declaretheoremstyle[
	bodyfont=\normalfont,
	shaded={
		margin=1em,
		bgcolor={rgb}{0.9,0.9,0.9}
}]{mydfn}

\declaretheoremstyle[
bodyfont=\normalfont,
spacebelow=1em,
spaceabove=1em,
]{myej}

\declaretheoremstyle[
bodyfont=\normalfont,
spacebelow=1em,
spaceabove=1em,
shaded={
    margin=8pt,
    bgcolor=White,
    rulecolor=coolblack,
    rulewidth=1pt
}]{thej}

\declaretheoremstyle[
bodyfont=\normalfont,
spacebelow=1em,
spaceabove=1em,
thmbox=M
]{myeg}

% DEFINICIONES DE ENTORNOS DE TEOREMAS
\declaretheorem[
	name=Teorema,
	refname={teorema,teoremas},
	Refname={Teorema,Teoremas},
	style=mythm
]{thm}

\declaretheorem[
	name=Corolario,
	refname={corolario,corolarios},
	Refname={Corolario,Corolarios},
	style=myej
]{cor}

\declaretheorem[
	name=Proposici\'{o}n,
	refname={proposici\'{o}n,proposiciones},
	Refname={Proposici\'{o}n, Proposiciones},
	sharenumber=thm,
	style=myej
]{pro}

\declaretheorem[
name=Lema,
refname={lema,lemas},
Refname={Lema, Lemas},
sharenumber=thm,
style=myej
]{lm}

%\newtheorem*{cor}{Corolario}
\newtheorem*{lem}{Lema}

\declaretheorem[
	name=Definici\'{o}n,
	refname={definici\'{o}n,definiciones},
	Refname={Definici\'{o}n,Definiciones},
	style=mydfn
]{dfn}

\declaretheorem[
name=Ejercicio,
refname={ejercicio,ejercicios},
Refname={Ejercicio,Ejercicios},
style=myej,
numbered=no
]{ex}

\declaretheorem[
name=Ejercicio propuesto,
refname={ejercicio,ejercicios},
Refname={Ejercicio,Ejercicios},
style=thej,
]{th_ex}

\declaretheorem[
name=Ejemplo,
refname={ejemplo,ejemplos},
Refname={Ejemplo,Ejemplos},
style=myeg
]{eg}

\declaretheorem[
name=Observaci\'{o}n,
refname={observaci\'{o}n,observaciones},
Refname={Observaci\'{o}n,Observaciones},
style=myej,
numbered=no
]{obs}

\makeglossaries
\newglossaryentry{di}{name={DI},description={Dominio de integridad (ver \autoref{dfn:dominiointegridad})}}
\newglossaryentry{dip}{name={DIP},description={Dominio de ideales principales (ver \autoref{dfn:dominioidealesprincipales})}}
\newglossaryentry{dfu}{name={DFU},description={Dominio de factorización única (ver \autoref{dfn:dfu})}}
\newglossaryentry{de}{name={DE},description={Dominio euclídeo (ver \autoref{dfn:de})}}


%\setcounter{chapter}{-1}
%COMANDOS ÚTILES PARA DERIVADAS
\AtBeginDocument{\renewcommand{\d}{\mathrm{d}}}
\newcommand{\Dd}[1]{\frac{\mathrm{d}}{\mathrm{d}#1}}
\newcommand{\dd}[2]{\frac{\mathrm{d}#1}{\mathrm{d}#2}}
\newcommand{\mbf}[1]{\mathbf{#1}}
% COMANDOS ÚTILES PARA LA TEORÍA DE GRUPOS
%\newcommand{\normsub}{\mathbin{\triangleleft}}
\newcommand{\normsub}{\lhd}
\newcommand{\uds}[1]{\mathcal{U}(#1)}
\newcommand{\inv}[1]{#1^{-1}}
\newcommand{\ima}{\text{Im }}
\newcommand{\isom}{\simeq}
\newcommand{\autom}[1]{\text{Aut}(#1)}
\newcommand{\gen}[1]{\langle#1\rangle}
\newcommand{\biy}[1]{\text{Biy}(#1)}
\newcommand{\quo}[2]{{\left.\raisebox{.2em}{$#1$}\middle/\raisebox{-.2em}{$#2$}\right.}}
\newcommand{\appad}[1]{\mathrel{\mkern#1}}
\newcommand{\setdef}[1]{\left\{#1\right\}}
\newcommand{\car}[1]{\mathrm{car}(#1)}
%\DeclareMathSymbol{\varprod}{\mathop}{largesymbolsA}{16}

\newcommand{\hr}{\rule{\textwidth}{.4pt}}
\newcommand{\N}{\mathbb{N}}
\newcommand{\Z}{\mathbb{Z}}
\newcommand{\R}{\mathbb{R}}
\newcommand{\Q}{\mathbb{Q}}
\newcommand{\C}{\mathbb{C}}
\newcommand{\A}{\mathbb{A}}
\newcommand{\F}{\mathbb{F}}
\newcommand{\J}{\mathbb{J}}
\newcommand{\B}{\mathbb{B}}
\DeclareMathOperator{\diag}{diag}
\DeclareMathOperator{\traza}{traza}
\renewcommand{\P}{\mathcal{P}}
\renewcommand{\S}{\mathbb{S}}
\newcommand{\ZnZ}{\quo{\Z}{n\Z}}
\newcommand{\ZmZ}{\quo{\Z}{m\Z}}
\newcommand{\0}{\mathbf{0}}
\newcommand{\1}{\mathbf{1}}
\newcommand{\gr}{\text{gr }}
\newcommand\ninf[1]{\left|\left|#1\right|\right|_\infty}

%%%% MATRICES
\newcommand\mx[1]{\begin{bmatrix*}#1\end{bmatrix*}}

\renewcommand\qedsymbol{$\diamondsuit$}

\title{Apuntes de Ecuaciones Diferenciales}
\author{Rafael S\'{a}nchez S\'{a}nchez}
\begin{document}
\begin{titlepage}
	\includepdf{portada.pdf}
\end{titlepage}


Revisión del \today $ $ a las \currenttime.

\begin{center}

\end{center}

\tableofcontents
\part{Primer parcial}
% !TeX root = ../galois.tex

\chapter{Anillos, polinomios y cuerpos}

\section{Anillos}

A lo largo de este curso se supondrán conocidos los contenidos de la asignatura \textit{Estructuras Algebraicas}, se pueden encontrar unos apuntes de los mismos en: \url{https://github.com/knifecake/apuntes/raw/master/ea/apuntes-ea.pdf}.

\begin{dfn}[Anillo]
	Un \textbf{anillo} es una terna $(A, +, \cdot)$ donde $+: A \times A \to A$ es una operación a la que llamamos suma, $\cdot: A \times A \to A$ es otra operación a la que llamamos producto y se verifican las siguientes propiedades
	\begin{enumerate}
		\item El par $(A, +)$ es un grupo abeliano
		\item El producto $\cdot$ es asociativo
		\item Se cumplen las propiedades distributivas:
		\begin{align}
			\forall a, b , c \in A,\ a\cdot (b + c) = a\cdot b + a \cdot c \\
			\forall a, b , c \in A,\ (a + b) \cdot c = a\cdot c + b \cdot c
		\end{align}
	\end{enumerate}
\end{dfn}

Con la operación $+$ tenemos las siguientes propiedades
\begin{enumerate}
	\item Asociatividad: $(a+b)+c = a+(b+c)$
	\item Elemento neutro aditivo: $\exists! \0 \in A \mid \0+a = a$
	\item Elemento inverso aditivo: $\forall a \in A, \exists -a \in A \mid a + (-a) = \0$
	\item Conmutatividad aditiva: $\forall a, b \in A,\ a+b = b+a$
\end{enumerate}

Con la operación $\cdot$ tenemos las siguientes propiedades
\begin{enumerate}
	\item Asociatividad: $a\cdot (b \cdot c) = (a \cdot b) \cdot c$
	\item No siempre existe el neutro multiplicativo: $\1 \in A \mid a\cdot 1 = 1 \cdot a = a$
	\item No siempre el producto es conmutativo.
	\item No siempre existe inverso multiplicativo: $\inv{a} \mid a\cdot \inv{a} = \1$
	\item No siembre se da la conmutatividad multiplicativa: $a \cdot b = b\cdot a$
\end{enumerate}

\begin{pro}[Producto con 0 en anillos]
	$\forall a \in A,\ a\cdot \0 = \0$
\end{pro}

\begin{proof}
	$a \cdot \0 = a \cdot(\0 + \0) = a\cdot \0 + a\cdot \0 \implies \0 = a\cdot \0$
\end{proof}

Además, a lo largo de este curso vamos a referirnos únicamente a los anillos conmutativos con unidad (o unitario), que cumplen las siguientes definiciones.

\begin{dfn}[Anillo con unidad o anillo unitario]
	Sea $(A, +, \cdot)$ un anillo. Decimos que es un anillo con unidad o un \textbf{anillo unitario} si tiene elemento neutro multiplicativo, es decir, si $\exists \1 \in A \mid \forall a \in A, \1 a = a \1 = a$.
\end{dfn}

\begin{dfn}[Anillo conmutativo]
    Sea $(A, +, \cdot)$ un anillo. Decimos que es un \textbf{anillo conmutativo} si se cumple que:
    $$
        r\cdot s = s\cdot r,\ \forall r,s \in A
    $$
\end{dfn}

\begin{eg}[Ejemplos de anillos]
    $\Z$, $\Q$, $\R$ y $\C$ con la suma y producto usual verifican todas las definiciones de anillo, anillo conmutativo y anillo unitario.
\end{eg}

Vamos a considerar además el concepto de anillo de polinomios:
\begin{dfn}[Anillo de polinomios]
    Sea $R$ un anillo, definimos el \textbf{anillo de polinomios} $R[x]$ como:
    $$
        R[x] = \left\{ \sum_{i > 0}^{n} a_i \cdot x^i \mid a_i \in R,\ n \in \N \right\}
    $$
    Es fácil ver que $R[x]$ es un anillo pues la suma y el producto son transitivas y asociativas.
\end{dfn}
\begin{obs}
    Vamos a considerar algunas definiciones y convenciones menores.\\
    \begin{enumerate}
        \item Sea $p \in R[x]$, $p$ es un polinomio y escribimos:
        $$
            p(x) = a_0 + a_1 x + \ldots + a_n x^n
        $$
        Donde llamamos \textit{coeficientes} del polinomio a los $a_i$.
        \item Sea $p \in R[x] = \sum_{i > 0} a_i x^i$, denominamos grado de $p$ a:
        $$
            \delta(p) = \max \left\{ i \mid a_i \neq 0 \right\}
        $$
        \item Sea $p \in R[x] = a_0 + a_1 x + \ldots + a_n x^n$, llamamos \textit{coeficiente director} al coeficiente del término de mayor grado ($a_n$).
        \item  Sea $p \in R[x] = a_0 + a_1 x + \ldots + a_n x^n$, llamamos \textit{termino independiente} al coeficiente libre ($a_0$).
        \item Sea $p \in R[x]$ con todos los coeficientes nulos, entonces $p$ es el \textit{polinomio cero}.
        $$
            0 = \sum_{i > 0} 0 \cdot x^n
        $$
        Por convención, $\delta(0) = -\infty$.
    \end{enumerate}
\end{obs}

\begin{dfn}[Polinomio mónico]
    Sea $R[x]$ un anillo de polinomios, decimos que $p \in R[x]$ es \textbf{mónico} si y sólo si su \textit{término director} es $1$.
\end{dfn}

\begin{dfn}[Divisor de cero]
    Sea $R$ un anillo, decimos que $r \in R$ es un \textbf{divisor de cero} si satisface:
    $$
        \exists s \in R,\ s\neq 0\ :\ r\cdot s = \0
    $$
\end{dfn}
\begin{dfn}[Unidad de un anillo]
    Sea $R$ un anillo, decimos que $r \in R$ es una \textbf{unidad} si satisface:
    $$
        \exists s \in R\ :\ r\cdot s = \1
    $$
    Decimos entonces que $r \in \uds{R}$, con $\uds{R} = \left\{a \mid a\text{ es una unidad}\right\}$
\end{dfn}

\begin{dfn}[Dominio de integridad]
    Sea $R$ un anillo, $R$ es un \textbf{dominio de integridad} si no tiene divisores de $\0$.
\end{dfn}

\section{Cuerpos}

\begin{dfn}[Cuerpo]
    Sea $(A, +, \cdot)$ un anillo conmutativo con unidad. Diremos que $A$ es un cuerpo si $A^{\times} = A\setminus \{\0\}$ es cerrado por la segunda operación (el \textit{producto}).
\end{dfn}

\begin{obs} $ $
    \begin{itemize}
        \item $R$ es un cuerpo si $\uds{R} = R$.
        \item $\1 \in \uds{R}$, para todo $R$ anillo unitario.
    \end{itemize}
\end{obs}

\begin{pro}[Cuerpo y dominio de integridad]\label{pro:cuerpoDI}
    Sea $R$ un cuerpo, entonces $R$ es un dominio de integridad.
\end{pro}
\begin{proof}
    Vamos a ver que $R$ no tiene divisores de $\0$. Sea $r \in R^{\times} = R\setminus \{\0\}$, supongamos $\exists s \in R^{\times}$ tal que:
    $$
        r \cdot s = 0
    $$
    Como $r \in \uds{R} = R^{\times}$ pues $R$ es un cuerpo, entonces, $\exists t \in R$ tal que $t\cdot r = r\cdot t = \1$. Por tanto:
    $$
        \0 = t \cdot (r \cdot s) = (t \cdot r) \cdot s = \1 \cdot s = s
    $$
    Y $s = \0$ contradice la hipótesis. Concluimos con que $\not \exists r, s \in R$ tal que $r \cdot s = \0$
\end{proof}

\begin{pro}[Dominio de integridad en anillos de polinomios]
    Sea $R$ un anillo. Si $R$ es un dominio de integridad, entonces $R[x]$ es un dominio de integridad.
\end{pro}
\begin{proof}
    Sean $f, g \in R[x]^{\times}$, y $a_m, b_k$ sus términos directores respectivamente. Como $R$ es un dominio de integridad, $a_m \cdot b_k \neq \0$, que coincide con el término director de $f \cdot g$ y no es nulo. Por tanto, $R[x]$ es un dominio de integridad.
\end{proof}

\begin{pro}[Propiedad de cuerpo en anillos de polinomios]
    $R[x]$ nunca es un cuerpo.
\end{pro}
\begin{proof}
    Solo hay que comprobar que aunque $f(x) = x \in R[x]$, $f(x) \notin \uds(R[x])$. Y por tanto $\uds(R[x]) \neq R[x]$, lo que nos dice que $R[x]$ no es un cuerpo.
\end{proof}

\begin{pro}[Unidades en anillos de polinomios]
    Sea $R$ un anillo, si $R$ es un dominio de integridad, entonces $\uds{R} = \uds{R[x]}$.
\end{pro}

\begin{obs}
    Podemos definir anillos como \textit{extensión} de otros, al igual que hicimos con los anillos de polinomios:\\
    \begin{itemize}
        \item $\Z[\sqrt{d}] = \left\{a + b\sqrt{d} \mid a,b \in \Z \right\}$, con $d \neq e^2,\ \forall e \in \Z$ es un anillo y un dominio de integridad, pero no es un cuerpo.
        \item $\Q[\sqrt{d}] = \left\{a + b\sqrt{d} \mid a,b \in \Q \right\}$, con $d \neq e^2,\ \forall e \in \Z$ es un cuerpo. Decimos que $\left\{ 1, \sqrt{d} \right\}$ es una $\Q$-base de $\Q[\sqrt{d}]$, pues todos los elementos de $\Q[\sqrt{d}]$ se pueden expresar como combinación lineal de los elementos de la $\Q$-base.
    \end{itemize}
\end{obs}

\begin{dfn}[Subanillo]
    Sea $R$ un anillo, $S \subseteq R$, $\1 \in S$. Decimos que $S$ es un \textbf{subanillo} si:
    \begin{itemize}
        \item $S$ es cerrado por suma y producto.
        \item Todo elemento tiene opuesto, es decir, $\forall a \in S, \exists b \in S\ :\ a + b = \0$.
    \end{itemize}
\end{dfn}

\begin{dfn}[Subcuerpo]
    Sean $R$ un cuerpo, $S \subseteq R$. Decimos que $S$ es un \textbf{subcuerpo} si:
    \begin{itemize}
        \item $S$ es un subanillo de $R$
        \item Todo elemento no nulo tiene inverso, es decir, $\forall a \in S^{\times}, \exists b \in S^{\times}\ :\ a \cdot b = \1$
    \end{itemize}
\end{dfn}

\begin{eg}[Ejemplos de subanillos y subcuerpos]$ $
    \begin{itemize}
        \item $\Z$ es subanillo de $\Q$
        \item $\Q$ es subcuerpo de $\R$ y $\C$
        \item $\Z[\sqrt{d}]$ es subanillo de $\Q[\sqrt{d}]$
    \end{itemize}
\end{eg}

\section{Ideales}

\begin{dfn}[Ideal]
    Sea $R$ un anillo, e $I \subseteq S$. $I$ es un \textbf{ideal} si:
    \begin{enumerate}
        \item $\forall a, b \in I,\ a - b \in I$
        \item $\forall r \in R,\ \forall a \in I$ se satisface: $r\cdot a \in I$
    \end{enumerate}
    Los ideales triviales son $\{\0\}$ y $R$.
\end{dfn}

\begin{obs}
    Sea $R$ un anillo, denotamos al ideal generado por $a \in R$ como:
    $$
        \gen{a}
    $$
\end{obs}

\begin{pro}[Ideal propio]
    Sea $R$ un anillo, $I$ un ideal:
    $$
        I \subsetneq R \iff \1 \in I \iff I \cap \uds{R} \neq \varnothing
    $$
\end{pro}

\begin{obs}
    Sea $R$ un anillo, $I \leq R$ un ideal:
    $$
        I \leq R \iff I \cap \uds{R} = \varnothing
    $$
    $$
        I = R \iff I \cap \uds{R} \neq \varnothing
    $$
\end{obs}


\begin{pro}[Ideales y cuerpos]
    Sea $R$ un cuerpo, y sea $I$ un ideal de $R$ (escribimos $I \leq R$), entonces $I = \{\0\}$ o $I = R$, ($I$ es impropio).\\
    El recíproco también es cierto.
\end{pro}

\begin{proof}$ $
    \begin{itemize}
        \item $\left( \implies \right)$
        $$
            R \text{ cuerpo} \implies \uds{R} = R^{\times} \implies I = \uds{R} \cup \{\0\} \text{ o trivialmente } I = \{\0\}
        $$
        \item $\left( \Longleftarrow \right)$\\
        Sea $a \in R^{\times}$, $a \in \gen{a} \leq R$
        $$
            \{\0\} \neq I = \gen{a},\text{ entonces } I = R \implies \exists u \in I \cap \uds{R} \neq \varnothing \implies u \in \gen{a} \implies u = a\cdot r, \text{ con } r \in R
        $$
        y por tanto:
        $$
            1 = u \cdot u^{-1} = a \cdot r \cdot u^{-1} \implies a \in \uds{R} \implies R \text{ es un cuerpo}
        $$
    \end{itemize}
\end{proof}

\begin{eg}[Ejemplos de ideales]$ $
    \begin{enumerate}
        \item $n\Z \leq \Z$
        \item $I = \left\{ f \in \Z[x] \mid \text{el termino independiente de f es par} \right\}$
    \end{enumerate}
\end{eg}

\begin{dfn}[Ideal principal]
    Sea $R$ un anillo, $a \in R$ un elemento. El ideal generado por $a$:
    $$
        \gen{a} = \left\{a\cdot r \mid r \in R\right\} = aR
    $$
    se denomina \textbf{ideal principal} generado por $a$.
\end{dfn}

\begin{pro}[Propiedades de ideales] Sea $ R $ un anillo e $ I \leq R $ un ideal.
    \begin{enumerate}
        \item Sean $I, J \leq R$ ideales, entonces $I+J = \left\{a+b \mid a \in I,\ b \in J\right\} \leq R$ es un ideal.
        \item Sea $\mbf{a} \in \R^n$, entonces $I = \gen{\mbf{a}} = \left\{ a_1r_1 + \cdots + a_nr_n \mid r_i \in R \right\} \leq R$ es un ideal.
        \item $\quo{R}{I} = \left\{r + I \mid r \in R\right\}$ es un anillo.
        \item \label{pro:correspondencia}(Teorema de correspondencia) Existe una biyección de la forma:
        \begin{align*}
            \{J \leq R \mid I \subseteq J \subseteq R\} &\longrightarrow \{\quo{J}{I} \leq \quo{R}{I}\}\\
             J \appad{80mu} &\longmapsto \appad{4mu} \{r + I \mid r \in J\}
        \end{align*}
    \end{enumerate}
\end{pro}

\begin{obs}
    En particular, si en $R$ todo ideal es principal e $I \leq R$, en $\quo{R}{I}$ todo ideal es principal.
\end{obs}

\begin{ex}[H1.5]
    Sea $n$ un número natural. Prueba que $\Z_n = \ZnZ$ es un cuerpo si y sólo si $n$ es primo.\\\\
    \begin{itemize}
        \item $\left(\ \Longleftarrow\ \right)$\\\\
        $ n $ primo $ \implies \forall k:\ 0 < k < n \text{ se cumple que } \mathrm{mcd}(k, n) = 1$, y por Bezout:
        $$
            1 = ka + nb,\ \text{ con } a, b \in \Z
        $$
        Donde el término $nb \equiv 0$ en $\ZnZ$ y por tanto queda $1 = ka$, lo que quiere decir que $k$ es el inverso de $a$ en $\ZnZ$.
        \item $\left(\implies\right)$\\\\
        Partimos de que $\ZnZ$ es cuerpo, por la proposición \ref{pro:cuerpoDI} sabemos que $\ZnZ$ es un dominio de integridad. Supongamos $n$ no primo, entonces $n = a\cdot b$, entonces:
        $$
            n \equiv \0 (\mathrm{mod}\ n) \implies \0 = (a + n\Z)(b + n\Z)
        $$
        Pero es imposible, ya que $a$ y $b$ serían divisores de $\0$ pero estamos en un dominio de integridad. Por tanto, $n$ es necesariamente primo.
    \end{itemize}
\end{ex}

\begin{ex}[H1.12]
    ¿Cuántos elementos tiene el anillo $\quo{\Z[i]}{\gen{2i}}$?¿Se trata de un cuerpo?\\\\
    Comenzamos escribiendo los conjuntos que forman parte del cociente:
    $$
        \Z[i] = \setdef{a + bi \mid a, b \in \Z}
    $$
    $$
        \gen{2i} = \gen{2} = 2\Z[i] = \setdef{2(a+bi) \mid a,b \in \Z} = (2\Z)[i] = \setdef{a + bi \mid a, b \in 2\Z}
    $$
    El conjunto cociente es por tanto:
    $$
        \quo{\Z[i]}{\gen{2i}} = \quo{\Z[i]}{2\Z[i]} = \left\{ a+bi + 2\Z[i] \mid a, b \in \Z \right\}
    $$
    Donde se tiene que:
    $$
        a + bi + 2\Z[i] = a_1 + b_1i + 2\Z[i] \iff a-a_1 \in 2\Z \text{ y } b-b_1 \in 2\Z \iff \setdef{a+bi+2\Z[i] \mid a, b \in \setdef{0, 1}} = \setdef{0, 1, i, 1+i}
    $$
    De esta forma vemos que el anillo tiene $4$ elementos y además no es un cuerpo ya que $i$ no tiene inverso.
\end{ex}

\begin{dfn}[Ideal primo]
    Sea $R$ un anillo e $I \leq R$ un ideal, diremos que $I$ es un \textbf{ideal primo} si:
    $$
        a\cdot b \in I \implies a \in I \text{ ó } b \in I
    $$
\end{dfn}

\begin{dfn}[Ideal maximal]
    Sea $R$ un anillo e $I \leq R$ un ideal, diremos que $I$ es un \textbf{ideal maximal} si:
    $$
        I \subseteq J \leq R \implies J = I \text { ó } J = R
    $$
\end{dfn}
\begin{thm}[Cociente de ideales primos y maximales]
    Sea $R$ un anillo, $I \leq R$ un ideal:\\
    \begin{enumerate}
        \item $I$ es primo $\iff \quo{R}{I}$ es un dominio de integridad.
        \item $I$ es maximal $\iff \quo{R}{I}$ es un cuerpo.
        \item $I$ ideal maximal $\implies$ $I$ ideal primo.
    \end{enumerate}
\end{thm}
\begin{proof}$ $
    \begin{enumerate}
        \item Se deja como ejercicio. Es directa usando definiciones.
        \item $I$ es maximal $\iff \quo{R}{I}$ no tiene ideales propios (por el teorema de correspondencia \ref{pro:correspondencia}). Y ya sabemos que $\quo{R}{I}$ no tiene ideales propios $\iff \quo{R}{I}$ es un cuerpo.
        \item Se sique de los apartados anteriores junto a la proposición $\ref{pro:cuerpoDI}$ que nos dice que un cuerpo es un dominio de integridad.
    \end{enumerate}
\end{proof}

\section{Homomorfismos}

\begin{dfn}[Homomorfismo de anillos]
    Sean $R, S$ anillos, $\varphi: R \to S$ es un \textbf{homomorfismo de anillos} si:\\
    \begin{enumerate}
        \item $\varphi$ es homomorfismo de grupos, es decir, $\varphi(0) = 0$ y $\varphi(a - b) = \varphi(a) - \varphi(b)$.
        \item $\varphi(1) = 1$.
        \item $\varphi(ab) = \varphi(a)\varphi(b)$.
    \end{enumerate}
\end{dfn}

\begin{obs}$ $
    \begin{itemize}
        \item $\ker{\varphi} = \setdef{a \in R \mid \varphi(a) = 0} \leq R$.
        \item $\varphi(R) \subseteq S $ es un subanillo. (No es ideal en general).
        \item $\varphi$ sobreyectiva, es decir, $\varphi$ es un epimorfismo $\iff \varphi(R) = S$.
    \end{itemize}
\end{obs}

\begin{obs}
    Si $R$ y $S$ son cuerpos y $\varphi: R \to S$ es un homomorfismo de anillos, llamaremos a $\varphi$ homomorfismo de cuerpos. Además $\varphi$ es inyectivo pues:
    $$
        1 \notin \ker{\varphi} \leq R \text{ cuerpo} \implies \ker{\varphi} = 0
    $$
\end{obs}

\begin{eg}[Proyección canónica]
    Sea $R$ un anillo, $I \leq R$ un ideal, es fácil ver que $\pi: R \to \quo{R}{I};\ r \mapsto r + I$ es un epimorfismo de anillos con $\ker{\pi} = I$.
\end{eg}

\begin{obs}
    $$\quo{R}{\ker{\varphi}} = \quo{R}{I}$$
\end{obs}

\begin{thm}[Teorema de isomorfía]
    Sea $\varphi: R \to S$ un homomorfismo de anillos, se tiene que:

    \begin{align*}
        \bar{\varphi}: \quo{R}{\ker \varphi} &\longrightarrow \varphi(S)\\
        r + \ker \varphi &\longmapsto \bar{\varphi}(r + \ker \varphi) = \varphi(r)
    \end{align*}
    es un isomorfismo de anillos.
\end{thm}

\begin{proof}
    Se deja como ejercicio.
\end{proof}

\begin{obs}
    Sea $\pi$ la proyección canónica, $\bar{\pi} = id_{\quo{R}{I}}$
\end{obs}


\begin{ex}[H1.14]
    Demuestra que si $\varphi: R \to S$ es un homomorfismo de anillos y $a \in \uds{R}$, entonces $\varphi(a) \in \uds(S)$. ¿Es cierto el recíproco?.\\\\
    Si $a \in \uds{R}$, tntonces $\exists b \in R$ tal que $\1 = a\cdot b$. Por tanto:
    $$
        \1 = \varphi(\1) = \varphi(a \cdot b) = \varphi(a) \cdot \varphi(b) \implies \varphi(a) \in \uds(S)
    $$
    El recíproco solo es cierto si $\varphi$ es un isomorfismo, pero en general no. Como contraejemplo consideramos el homomorfismo identidad $\iota : \Z \to \Q;\ a \mapsto a$. Es fácil ver que es un homomorfismo de anillos, sin embargo: $\iota(2) = (2)$ pero $\iota(2) \in \uds{\Q}$ y $2 \notin \uds{\Z}$.
\end{ex}

\begin{ex}[H1.16]
    Demuestra que:
    \begin{enumerate}
        \item No existe ningún homomorfismo de anillos $\varphi: \Q \to \Z_p$ para $p \in \Z$ primo.
        \item No existe ningún homomorfismo de anillos $\varphi: \R \to \Q$.
    \end{enumerate}
    Solución:
    \begin{enumerate}
        \item Sea $\varphi: \Q \to \Z_p;\ \1 \mapsto \1 + p\Z$.
        $$
            \varphi(p) = \varphi\left(\sum_1^p 1 \right) = \sum_1^p (\1 + p\Z) = p + p\Z = 0.
        $$
        y como $p \in \uds{\Q}$, es imposible que la imagen de una unidad no sea otra por medio de un homomorfismo, por tanto, dicho homomorfismo no existe.
        \item Sea $\varphi: \R \to \Q;\ \sqrt{2} \mapsto a$
        $$
            2 = \varphi(1 + 1) = \varphi(2) = \varphi(\sqrt{2}^2) = \varphi(\sqrt{2})^2 = a^2,\ a \in \Q
        $$
        que es una contradicción pues no existe dicho $a$, con lo que no existe el homomorfismo.
    \end{enumerate}
\end{ex}

\begin{ex}[H1.21]
    Fijado un entero $n \in \Z$ con $n \geq 2$, demuestra que el anillo cociente $\quo{\Z[x]}{n\Z[x]}$ es isomorfo a $\Z_n[x]$. Conclute que el ideal $n\Z[x]$ es primo si y sólo si $n$ es un número primo.\\\\
    Vamos a dar una guía de como proceder con el ejercicio:
    $$
        \text{Sea } \varphi: \Z[x] \to \Z_n[x];\ (a_0 + \ldots + a_n x^n) \mapsto (\bar{a_0} + \ldots + \bar{a_n} x^n)
    $$
    donde $\bar{a_i} = a_i + n\Z \in \quo{\Z}{n\Z}$.
    \begin{itemize}
        \item Comprobar que $\varphi$ es un homomorfismo de anillos.
        \item Comprobar que $\varphi$ es sobreyectiva.
        \item Ver que $\ker \varphi = n\Z[x]$.
        \item Aplicar el teorema de isomorfía.
    \end{itemize}
\end{ex}

\begin{eg}[Homomorfismo de evaluación]
    Sea $R$ un anillo, $a \in R$.
    \begin{align*}
        \mathcal{E}_a : R[x] &\longleftarrow R\\
                     f(x) &\longmapsto f(a)
    \end{align*}
    es un homomorfismo de anillos sobreyectivo.
\end{eg}

\begin{obs}
    Si $R = K$ es un cuerpo:
    $$
        \quo{K[x]}{\ker \mathcal{E}_a} \isom K \implies \ker\mathcal{E}_a \text{ es maximal.}
    $$
\end{obs}

\section{Anillos de polinomios}

\begin{pro}[Algoritmo de la división]
    Sea $R$ un anillo, $f, g \in R[x]^{\times}$ polinomios con coeficientes en $R$. Si el coeficiente director de $g$ es una unidad de $R$, entonces $\exists d,r \in R[x]$ únicos tales que:
    $$
        f = g \cdot d + r \text{ con } \delta(r) < \delta(g)
    $$
    Diremos que $g \divides f$ si $r = \0$.
\end{pro}
\begin{dfn}[Raíz de un polinomio]
    Sea $R$ un anillo, $f \in R[x]^{\times}$ un polinomio, decimos que $a \in R$ es una \textbf{raíz} de $f$ si $\mathcal{E}_a(f) = f(a) = \0$
\end{dfn}
\begin{cor}[Ruffini]\label{cor:ruffini}
    Sea $R$ un anillo, $f \in R[x]^{\times}$ un polinomio:
    $$
        a \text{ es raíz de } f \iff f(x) = (x-a)\cdot g(x)
    $$
\end{cor}
\begin{proof}$ $
    \begin{itemize}
        \item $\left(\  \Longleftarrow\ \right)$
        $$
            \mathcal{E}_a(f) = \mathcal{E}_a(x-a) \cdot \mathcal{E}_a(g) = \0
        $$
        \item $(\implies)$
        $$
            f(x) = (x - a)\cdot d(x) + r(x);\ \delta(r) \leq \delta(x-a) \implies r \in R;\ f(a) = r = 0 \implies g(x) = d(x)
        $$
    \end{itemize}
\end{proof}

\begin{eg}[Uso de Ruffini]
    Sea $f(x) = x^2 + x + 1$, $f(x) \in \Z_3[x]$.\\ Es fácil ver que $f(1) = 0$, según Ruffini (corolario \ref{cor:ruffini}) $(x-1) \divides f$. Y es cierto, de hecho: $f(x) = (x - 1)(x - 1)$.
\end{eg}

\begin{thm}[Raíces y dominio de integridad]
    Sea $R$ un dominio de integridad, $f \in R[x]^{\times}$ un polinomio y $\alpha_1, \ldots, \alpha_n$ raíces distintas de $f$, entonces $n \leq \delta(f)$.
\end{thm}


\part{Apéndices}

%\renewcommand{\chaptername}{Apéndice}
%\renewcommand{\thechapter}{\Alph{chapter}}
%\setcounter{chapter}{0}

%% !TeX root = ../apuntes-ea.tex

\chapter{Ejercicios}

\section{Hoja 1}

\begin{ex}[H1.2]
	Sean $a,b,c \in G = (-1,1)$. Probamos las propiedades de los grupos.
	\begin{itemize}
		\item \textbf{Asociatividad:}
		\begin{align*}
			(a \ast b) \ast c = \left(\frac{a+b}{1+ab}\right) \ast c = \frac{\left(\frac{a+b}{1+ab}\right) + c}{1 + \left(\frac{a+b}{1+ab}\right)c} = \frac{\frac{a+b+c+abc}{1+ab}}{\frac{1+ab+ac+bc}{1+ab}} = \frac{a+b+c+abc}{1+ab+ac+bc} \\
			a \ast (b \ast c) = a \ast \left(\frac{b+c}{1+bc}\right) = \frac{a+\left(\frac{b+c}{1+bc}\right)}{1+a\left(\frac{b+c}{1+bc}\right)} = \frac{\frac{a + b + c + abc}{1+bc}}{\frac{1+ab+ac+bc}{1+bc}} = \frac{a+b+c+abc}{1+ab+ac+bc}
		\end{align*}
		\item \textbf{Elemento neutro:} es el $0$ ya que $x\ast 0 = \frac{x+0}{1+x\cdot 0} = \frac{x}{1} = x$ y además $0 \ast x = \frac{0+x}{1+0\cdot x} = \frac{x}{1} = x$
		\item \textbf{Elemento inverso:} la ecuación
		\begin{align*}
			x \ast \inv{x} = 0 \iff \frac{x + \inv{x}}{1+x\inv{x}} = 0 \iff \inv{x} = -x
		\end{align*}
		siempre tiene solución y ocurre lo mismo para la ecuación $\inv{x} \ast x = 0 \iff \inv{x} = -x$
		\item \textbf{Clausura:} tenemos que probar que si $x,y \in (-1,1)$ entonces $x \ast y \in (-1,1)$. Consideramos $f(x,y) = x \ast y = \frac{x+y}{1+xy}$. Derivando tenemos que $\nabla f(x,y) = (\frac{1}{(1+xy)^2},\ \frac{1}{(1+xy)^2}) \neq 0, \forall x,y \in [-1,1]\times[-1,1]$. Si el máximo no se alcanza en ningún sitio de dentro del cuadrado $(-1,1)\times(-1,1)$ se tendrá que alcanzar en el borde.
		\begin{itemize}
			\item Fijado $x = 1$ tenemos que $f(1,y) = \frac{1+y}{1+y} = 1 \implies f(1, -1 < y < 1) < 1$ porque si $f(1, -1 < y < 1)$ tomara un valor mayor que $1$ habría un máximo en $(-1,1)\times(-1,1)$ y esto no puede ser pues $\nabla f$ no se anula en el cuadrado.
			\item Fijado $x = -1$ tenemos que $f(-1, y) = \frac{y-1}{1-y} = -1 \implies f(-1, -1 < y < 1) > -1$ por la misma razón que antes.
			\item Hacemos lo mismo fijando la $y$ y variando la $x$.
		\end{itemize}
		En el borde (que no está incluido) se alcanzan máximo y mínimo que acotan a $f$ en el cuadrado:
		\begin{align*}
			-1 < f(x,y) = x \ast y < 1,\quad \forall x, y \in G
		\end{align*}
	\end{itemize}
\end{ex}

\begin{ex}[H1.3]Hallar los inversos de los siguientes elementos, cada uno en su grupo correspondiente:
	\begin{enumerate}
		\item $o(\overline{11})$ en $\uds{\Z^\ast/23\Z}$ es $22$ porque $23 \cdot 22 \equiv 1 \mod 23$
		\item $o(\overline{5})$ en $\uds{\Z^\ast/31\Z}$ es $3$ porque $5 \cdot 3 \equiv 1 \mod 31$
	\end{enumerate}
\end{ex}

\begin{ex}[H1.33]
	\label{ex:h1.33}
	Sea $G$ un grupo. Suponed que existe un único $a \in G$ de orden 2. Demostrad que $a \in Z(G)$.
\end{ex}

\begin{proof}
	Recordamos que $a \in Z(G) \iff ga = ag,\ \forall g \in G$. Definimos el isomorfismo de conjugación $\phi_g (x) = gx\inv{g}$ para algún $g$. Como $\phi_g$ es isomorfismo lleva elementos de orden $n$ en elementos de orden $n$. Entonces $\phi_g(a) = a$ ya que $a$ es el único elemento de orden 2. Por tanto $ga\inv{g} = a \implies ga = ag \implies a \in Z(G)$.
\end{proof}

\section{Hoja 2}

\begin{ex}[H2.1]
	\label{ex:h2.1}
	Se considera el tercer grupo diédrico $D_3$. Se pide hallar lo siguiente:
	\begin{enumerate}
		\item Las clases de conjugación de cada uno de sus elementos.
		\begin{proof}
			Las clases dan una partición del grupo. Si un elemento pertenece a una clase, entonces la clase de ese elemento también es la clase a la que pertenece.
			\begin{itemize}
				\item $cl(e) = \{e\}$
				\item $cl(B)$? Sabemos que $|cl(B)| = [G:C(B)]$. Sabemos que $\gen{B} =\{1, B, B^2\} \subset C(B)$ luego $|C(B)| \geq 3$. Si hubiera más elementos en $C(B)$ tendríamos que $|C(B)| = 6$ pues $C(B) < D_3$. Esto no ocurre porque sabemos que $B$ no conmuta con todos los demás elementos. Por ejemplo $BA \neq AB$. Por tanto $|C(B)| = 3 \implies |cl(B)| = [D_3:C(B)]= 6/3 = 2$. Es claro que $B \in cl(B)$. Además, como $cl(B)$ contiene elementos transformados por el isomorfismo conjugación sabemos que el otro elemento que hay tiene orden 3. El único elemento que queda de orden 3 es $B^2 \implies cl(B) = \{B, B^2\}$.
				
				\item $cl(A)?$ Sabemos que $A$ no conmuta con todos ($A \not\in Z(D_3)$) luego $|C(A)| < 6$. Sabemos que $\gen{A} = \{1, A\} < C(A)$. Además, como $C(A)$ es un (sub)grupo sabemos que no puede haber más elementos porque si los hubiera, $|\gen{A}| \divides |C(A)| \implies C(A) \geq 6$ pero ya hemos visto que no puede ser. Es decir que $|cl(A)| = [D_3:D(A)] = 6 / 2 = 3$. Por tanto $cl(A)$ incluye los 3 elementos que nos quedan: $cl(A) = \{A, AB, AB^2\}$.
			\end{itemize}
		\end{proof}
	
		\item Los elementos de $\text{Int}(D_3)$.
		\item Los centralizadores $C_{D_3}(x)$ para cada $x \in D_3$
		\item Los normalizadores $N(H)$ para cada $H < D_3$.
	\end{enumerate}
\end{ex}

\begin{ex}[H2.2]
	
	\begin{proof}
		Obtenidas las clases en el ejercicio \nameref{ex:h2.1} se verifica que $|D_3| = |cl(e)| + |cl(B)| + |cl(A)| = 1 + 2 + 3 = 6$
	\end{proof}
	
\end{ex}

\begin{ex}[H2.6]
	Sea $G$ un grupo. ¿Verdadero o falso?
	\begin{enumerate}
		\item $H < G$ y $H$ conmutativo implica $H \normsub G$.
		\item $H < G$ y $|H| = 2$ implica $H \normsub G$.
		\item Si $\varphi: G \to G_1$ es un homomorfismo de grupos, entonces $\ima \varphi \normsub G$
		\item Si $H \normsub K$ y $K \normsub G$ entonces $H \normsub G$
		\item Si $H \normsub G$ y $|H| = m$ entonces $H$ es el único subgrupo de $G$ de orden $m$.
		\item Si $H \normsub G$ entonces $H < Z(G)$.
		\begin{proof}[FALSO]
			Contraejemplo: En $G = D_4$ tomamos $H = \gen{B^2} = \{1, B, B^2, B^3\} \not\subset Z(D_4) = \{1, B^2\}$.
		\end{proof}
	\end{enumerate}
\end{ex}

\begin{ex}[H2.10]
	\begin{proof}
		Fijado $n$ y definida $\alpha_n : G \to G,\ x \mapsto x^n$ tenemos que $\alpha_n$ es un homomorfismo de grupos. Además podemos expresar $H_2 = \ker \alpha_n \implies H_2 \normsub G$. Además también tenemos que $H_1 = \ima \alpha_n < G$. Veamos que $H_1 \normsub G$. Es decir, que $gH_1\inv{g} = H_1,\ \forall g \in G$. Para ello tomamos $x_1^n \in H_1$ y lo conjugamos $gx_1^n\inv{g} = (gx_1\inv{g})^n$ por ser $\alpha$ homomorfismo de grupos. En particular $(gx_1\inv{g})^n \in \ima \alpha \implies (gx_1\inv{g})^n \in H_1 \implies (gx_1\inv{g})^n = x_2^n$ para algún $x_2 \in H_1 \implies H_1 \normsub G$.
	\end{proof}
\end{ex}

\begin{ex}[H2.13] Si $A$ es un grupo abeliano con $n$ elementos y $k$ es un entero primo con $n$, demostrad que la aplicación $\varphi : A \to A$ definida por $\varphi(a) = a^k$ es un isomorfismo.
	\begin{itemize}
		\item $\varphi$ homomorfismo de grupos.
		\begin{proof}
			\begin{align*}
				\varphi(a)\varphi(b) = a^kb^k = (ab)^k = \varphi(ab)
			\end{align*}
		\end{proof}
		\item $\varphi$ biyectiva $\iff \varphi$ inyectiva ya que dominio y codominio coinciden
		\begin{proof}
			$\ker \varphi = \{a \in A \mid \varphi(a) = a^k = 1\}$. Probaremos que $a^k = 1 \iff a = 1$ y por tanto que $\ker \varphi = \{1\} \implies \varphi$ inyectiva. Sabemos que $a^k = 1 \iff o(a^k) = 1$. Sea $t = o(a) \divides n$. Distinguimos dos casos
			\begin{itemize}
				\item Si $t = 1$ entonces $a = 1$ y ya está
				\item Si $t > 1$ entonces $o(a^k) = \frac{t}{mcd(k, t)} = \frac{t}{1} > 1$ contradicción. Luego necesariamente $t = o(a) = 1$.
			\end{itemize}
		\end{proof}
	\end{itemize}
\end{ex}

\begin{ex}[H2.22]
	\label{ex:h2.22}
	Demostrad que si $G$ es un grupo no conmutativo y tiene orden $p^3$ ($p$ un número primo) entonces $Z(G)$ tiene orden $p$.
	
	\begin{proof}
		Sabemos que $Z(G) < G \implies |Z(G)| \divides |G| \implies |Z(G)| \in \{1, p, p^2, p^3\}$
		\begin{itemize}
			\item $|Z(G)| \neq p^3$ porque en tal caso $G$ sería conmutativo
			\item $|Z(G)| \neq 1$ porque $G$ es un p-grupo y por tanto su centro no es el trivial.
			\item Si $|Z(G)| = p^2$ entonces $|G/Z(G)| = p \implies G/Z(G)$ es cíclico lo que no es posible si $G$ no es abeliano.
		\end{itemize}
	
		Por descarte concluimos que $|Z(G)| = p$.
	\end{proof}
\end{ex}


\begin{ex}[H2.25]
	Sabemos que $\autom{\Z/12\Z} \isom \uds{\Z^\star/12\Z}$ donde $\Z^\star/12\Z$ es el grupo multiplicativo $(\{\overline{1}, \overline{2}, \overline{3}, \dots,  \overline{11}\}, \cdot)$. Queda $\uds{\Z^\star/12\Z} = (\{\overline{1}, \overline{5}, \overline{7}, \overline{11}\}, \cdot )$ y además da la casualidad que $\forall x \in \uds{\Z^\star/12\Z},\ o(x) = 2$ (todos los elementos son su propio inverso) por lo que no tenemos restricciones al definir $f: \Z/2\Z \to \autom{\Z/12\Z}$:
	\begin{align*}
		f:\Z/2\Z & \to \autom{\Z/12\Z} \isom \uds{\Z^\star/12\Z} \\
		e = \overline{0} &\mapsto 1 \\
		\overline{1} &\mapsto \{\overline{1}, \overline{5}, \overline{7}, \overline{11}\}
	\end{align*}
\end{ex}

\begin{ex}[H2.26]
	Sea $|G_1| = m, |G_2| = n,\ mcd(m, n) = 1$. Si $f:G_1 \to G_2$ es h. de g. sabemos que $o(f(a)) \divides o(a),\ \forall a \in G_1$. Además $o(a) \divides m \land o(f(a)) \divides n$ por el teorema de Lagrange (\ref{thm:lagrange}).
	\begin{align*}
		\begin{cases}
		o(a) \divides m \land o(f(a)) \divides n \\
		mcd(m, n) = 1 \\
		o(f(a)) \divides o(a)
		\end{cases} \implies o(a) = o(f(a)) = 1, \forall a \in G_1
	\end{align*}
	Por lo que solo puede haber un homomorfismo entre ellos y además es el trivial $f(a) = e_{G_2}$.
\end{ex}

\begin{ex}[H2.19]
	Definimos una función $f:[0, 2\pi] \subset \R \to \mathbb{S}^1, \alpha\mapsto \cos \alpha + i \sin \alpha$. Esta función tiene la propiedad de que $f(\alpha) \cdot f(\alpha') = \cos(\alpha + \alpha') + i\sin(\alpha + \alpha') = f(\alpha + \alpha')$ y por tanto es un h. de g.\footnote{A la izquierda (en $\R$) sumamos pero a la derecha (en $\mathbb{S}^1$) multiplicamos.} entre $R$ y $\mathbb{S}^1$.
	
	Un elemento de $\cos \alpha + i \sin \alpha \in \mathbb{S}^1$ es de torsión $\iff \exists n \mid (\cos \alpha + i \sin \alpha )^n = 1$. Ahora bien $(\cos \alpha + i \sin \alpha )^n = \cos n\alpha + i \sin n\alpha = 1 \iff n\alpha = k 2\pi$. 
\end{ex}

\section{Hoja 4}

\begin{ex}[H4.11] Hallar los subgrupos de Sylow de $S_5$.
	Sabemos que $|S_5| = 5! = 2^3\cdot3\cdot 5$ y por \nameref{thm:sylow1} tenemos lo siguiente:
	\begin{itemize}
		\item $\exists P_2,\ |P_2| = 2^3 = 8$.
		\item $\exists P_3,\ |P_3| = 3$. Además en $S_5$ hay $\binom{5}{3}2! =20 $ 3-ciclos y en cada $gP_3\inv{g} = \{1, a, a^2 \mid o(a) = 3\}$ hay 2 elementos de orden 3 distintos. Además, $g_1P_3\inv{g_1} \cap g_2P_3\inv{g_2} = \{e\}$ porque si su intersección fuera más grande entonces serían el mismo subgrupo (porque son cíclicos). Es por esto que tenemos que repartir 40 elementos dando 2 a cada 3-grupo con lo que obtenemos $n_3 = 20 / 2 = 10$ 3-subgrupos de Sylow en $S_5$.
		\item $\exists P_5,\ |P_5| = 5$. Además en $S_5$ hay $4! = 24$ 5-ciclos (elementos de orden 5) y en cada $gP_5\inv{g} = \{1, a, a^2, a^3, a^4 \mid o(a) = 5\}$ tenemos $4$ elementos de orden $5$ distintos. Además, $g_1P_5\inv{g_1} \cap g_2P_5\inv{g_2} = \{e\}$ porque si su intersección fuera más grande entonces serían el mismo. Así, tenemos $24$ 5-ciclos a repartir entre los diferentes $gP_5\inv{g}$ dando 4 5-ciclos a cada 1. Por tanto tenemos $n_5 = 24/4 = 6$ 5-subgrupos de Sylow en $S_5$.
	\end{itemize}
\end{ex}

\begin{ex}[H4.18]
	Demostrar que todo grupo de orden $|G| = 5^3 \cdot 7^3$ tiene un subgrupo normal de orden 125.
	
	\begin{proof}
		\nameref{thm:sylow1} $\implies \exists P_5 < G,\ |P_5| = 125$. \nameref{thm:sylow3} $\implies n_5 \divides 7^3 \land n_5 \equiv 1 \mod 5$ es decir $n_5 \in \{1, 7, 49, 343\} \land n_5 \in \{1, 6, 11, \dots\}$. Como ni $49$ ni $343$ son conguentes con $1$ módulo 5 tenemos que $n_5 = 1 \implies P_5 \normsub G$.
	\end{proof}
\end{ex}

\begin{ex}[H4.20] Hallad todos los grupos abelianos de órdenes 36, 64, 96 y 100.
	\begin{enumerate}
		\item $|G| = 36 = 2^2\cdot3^2$
		\begin{proof}
			\nameref{thm:sylow1} $\implies \exists P_2,\ |P_2| = 4 \land \exists P_3,\ |P_3| = 9$. Además $G$ abeliano $\implies P_2, P_3 \normsub G \implies G \isom P_2 \times P_3$. Estudiamos los grupos de orden 4 y de orden 9
			\begin{itemize}
				\item $|P_2| = 4$ entonces $P_2 \isom \Z/2\Z \times \Z/2\Z \lor P_2 \isom \Z/4\Z$
				\item $|P_3| = 9$ entonces $P_3 \isom \Z/3\Z \times \Z/3\Z \lor P_3 \isom \Z/9\Z$
			\end{itemize}
			Como $G \isom P_2 \times P_3$ tenemos 4 posibles grupos abelianos de orden 36.
		\end{proof}
	\end{enumerate}
\end{ex}

\begin{ex}[H4.22]
	Hallar todos los grupos abelianos de orden 175.
	
	\begin{proof}
		$|G| = 5^2\cdot 7$. Por el \nameref{thm:sylow1} tenemos que $\exists P_5, P_7 < G$ con $|P_5| = 25,\ |P_7| = 7$ y además por ser $G$ abeliano tenemos que $P_5, P_7 \normsub G \implies G \isom P_5 \times P_7$. Estudiamos los grupos de orden 25 y de orden 7:
		\begin{itemize}
			\item $|P_5| = 25 \land P_5$ abeliano $\implies P_5 \isom \Z/25\Z \lor P_5 \isom \Z/5\Z \times \Z/5\Z$. En ambos casos $P_5$ es producto directo de cíclicos pues $\Z/n\Z$ es cíclico.
			\item $|P_7| = 7 \land P_7$ abeliano $\implies P_7 \isom \Z/7\Z$. Ocurre lo mismo que con $P_5$.
		\end{itemize}
		Concluimos que $G \isom \Z/25\Z \times \Z/7\Z \lor G \isom \Z/5\Z \times \Z/5\Z \times \Z/7\Z$. Los dos casos son abelianos por ser producto directo de grupos cíclicos.
	\end{proof}
\end{ex}

\begin{ex}[H4.23]
	¿Cuántos elementos de orden 3 puede tener un grupo abeliano de orden 36?
	
	\begin{proof}
		$|G| = 36 = 2^2 3^2$. \nameref{thm:sylow1} $\implies \exists P_2, P_3 < G,\ |P_2| = 4,\ |P_3| = 9$. Estudiamos los grupos de órdenes 4 y 9:
		\begin{itemize}
			\item $|P_3| = 9 \implies P_3 \isom \Z/9\Z \lor P_3 \isom \Z/3\Z \times \Z/3\Z$
			% TODO: terminar
		\end{itemize}
	\end{proof}
\end{ex}

%\renewcommand{\chaptername}{Apéndice}
%\renewcommand{\thechapter}{\Alph{chapter}}
%\setcounter{chapter}{0}

%% !TeX root = ../apuntes-ea.tex

\chapter{Ejercicios}

\section{Hoja 1}

\begin{ex}[H1.2]
	Sean $a,b,c \in G = (-1,1)$. Probamos las propiedades de los grupos.
	\begin{itemize}
		\item \textbf{Asociatividad:}
		\begin{align*}
			(a \ast b) \ast c = \left(\frac{a+b}{1+ab}\right) \ast c = \frac{\left(\frac{a+b}{1+ab}\right) + c}{1 + \left(\frac{a+b}{1+ab}\right)c} = \frac{\frac{a+b+c+abc}{1+ab}}{\frac{1+ab+ac+bc}{1+ab}} = \frac{a+b+c+abc}{1+ab+ac+bc} \\
			a \ast (b \ast c) = a \ast \left(\frac{b+c}{1+bc}\right) = \frac{a+\left(\frac{b+c}{1+bc}\right)}{1+a\left(\frac{b+c}{1+bc}\right)} = \frac{\frac{a + b + c + abc}{1+bc}}{\frac{1+ab+ac+bc}{1+bc}} = \frac{a+b+c+abc}{1+ab+ac+bc}
		\end{align*}
		\item \textbf{Elemento neutro:} es el $0$ ya que $x\ast 0 = \frac{x+0}{1+x\cdot 0} = \frac{x}{1} = x$ y además $0 \ast x = \frac{0+x}{1+0\cdot x} = \frac{x}{1} = x$
		\item \textbf{Elemento inverso:} la ecuación
		\begin{align*}
			x \ast \inv{x} = 0 \iff \frac{x + \inv{x}}{1+x\inv{x}} = 0 \iff \inv{x} = -x
		\end{align*}
		siempre tiene solución y ocurre lo mismo para la ecuación $\inv{x} \ast x = 0 \iff \inv{x} = -x$
		\item \textbf{Clausura:} tenemos que probar que si $x,y \in (-1,1)$ entonces $x \ast y \in (-1,1)$. Consideramos $f(x,y) = x \ast y = \frac{x+y}{1+xy}$. Derivando tenemos que $\nabla f(x,y) = (\frac{1}{(1+xy)^2},\ \frac{1}{(1+xy)^2}) \neq 0, \forall x,y \in [-1,1]\times[-1,1]$. Si el máximo no se alcanza en ningún sitio de dentro del cuadrado $(-1,1)\times(-1,1)$ se tendrá que alcanzar en el borde.
		\begin{itemize}
			\item Fijado $x = 1$ tenemos que $f(1,y) = \frac{1+y}{1+y} = 1 \implies f(1, -1 < y < 1) < 1$ porque si $f(1, -1 < y < 1)$ tomara un valor mayor que $1$ habría un máximo en $(-1,1)\times(-1,1)$ y esto no puede ser pues $\nabla f$ no se anula en el cuadrado.
			\item Fijado $x = -1$ tenemos que $f(-1, y) = \frac{y-1}{1-y} = -1 \implies f(-1, -1 < y < 1) > -1$ por la misma razón que antes.
			\item Hacemos lo mismo fijando la $y$ y variando la $x$.
		\end{itemize}
		En el borde (que no está incluido) se alcanzan máximo y mínimo que acotan a $f$ en el cuadrado:
		\begin{align*}
			-1 < f(x,y) = x \ast y < 1,\quad \forall x, y \in G
		\end{align*}
	\end{itemize}
\end{ex}

\begin{ex}[H1.3]Hallar los inversos de los siguientes elementos, cada uno en su grupo correspondiente:
	\begin{enumerate}
		\item $o(\overline{11})$ en $\uds{\Z^\ast/23\Z}$ es $22$ porque $23 \cdot 22 \equiv 1 \mod 23$
		\item $o(\overline{5})$ en $\uds{\Z^\ast/31\Z}$ es $3$ porque $5 \cdot 3 \equiv 1 \mod 31$
	\end{enumerate}
\end{ex}

\begin{ex}[H1.33]
	\label{ex:h1.33}
	Sea $G$ un grupo. Suponed que existe un único $a \in G$ de orden 2. Demostrad que $a \in Z(G)$.
\end{ex}

\begin{proof}
	Recordamos que $a \in Z(G) \iff ga = ag,\ \forall g \in G$. Definimos el isomorfismo de conjugación $\phi_g (x) = gx\inv{g}$ para algún $g$. Como $\phi_g$ es isomorfismo lleva elementos de orden $n$ en elementos de orden $n$. Entonces $\phi_g(a) = a$ ya que $a$ es el único elemento de orden 2. Por tanto $ga\inv{g} = a \implies ga = ag \implies a \in Z(G)$.
\end{proof}

\section{Hoja 2}

\begin{ex}[H2.1]
	\label{ex:h2.1}
	Se considera el tercer grupo diédrico $D_3$. Se pide hallar lo siguiente:
	\begin{enumerate}
		\item Las clases de conjugación de cada uno de sus elementos.
		\begin{proof}
			Las clases dan una partición del grupo. Si un elemento pertenece a una clase, entonces la clase de ese elemento también es la clase a la que pertenece.
			\begin{itemize}
				\item $cl(e) = \{e\}$
				\item $cl(B)$? Sabemos que $|cl(B)| = [G:C(B)]$. Sabemos que $\gen{B} =\{1, B, B^2\} \subset C(B)$ luego $|C(B)| \geq 3$. Si hubiera más elementos en $C(B)$ tendríamos que $|C(B)| = 6$ pues $C(B) < D_3$. Esto no ocurre porque sabemos que $B$ no conmuta con todos los demás elementos. Por ejemplo $BA \neq AB$. Por tanto $|C(B)| = 3 \implies |cl(B)| = [D_3:C(B)]= 6/3 = 2$. Es claro que $B \in cl(B)$. Además, como $cl(B)$ contiene elementos transformados por el isomorfismo conjugación sabemos que el otro elemento que hay tiene orden 3. El único elemento que queda de orden 3 es $B^2 \implies cl(B) = \{B, B^2\}$.
				
				\item $cl(A)?$ Sabemos que $A$ no conmuta con todos ($A \not\in Z(D_3)$) luego $|C(A)| < 6$. Sabemos que $\gen{A} = \{1, A\} < C(A)$. Además, como $C(A)$ es un (sub)grupo sabemos que no puede haber más elementos porque si los hubiera, $|\gen{A}| \divides |C(A)| \implies C(A) \geq 6$ pero ya hemos visto que no puede ser. Es decir que $|cl(A)| = [D_3:D(A)] = 6 / 2 = 3$. Por tanto $cl(A)$ incluye los 3 elementos que nos quedan: $cl(A) = \{A, AB, AB^2\}$.
			\end{itemize}
		\end{proof}
	
		\item Los elementos de $\text{Int}(D_3)$.
		\item Los centralizadores $C_{D_3}(x)$ para cada $x \in D_3$
		\item Los normalizadores $N(H)$ para cada $H < D_3$.
	\end{enumerate}
\end{ex}

\begin{ex}[H2.2]
	
	\begin{proof}
		Obtenidas las clases en el ejercicio \nameref{ex:h2.1} se verifica que $|D_3| = |cl(e)| + |cl(B)| + |cl(A)| = 1 + 2 + 3 = 6$
	\end{proof}
	
\end{ex}

\begin{ex}[H2.6]
	Sea $G$ un grupo. ¿Verdadero o falso?
	\begin{enumerate}
		\item $H < G$ y $H$ conmutativo implica $H \normsub G$.
		\item $H < G$ y $|H| = 2$ implica $H \normsub G$.
		\item Si $\varphi: G \to G_1$ es un homomorfismo de grupos, entonces $\ima \varphi \normsub G$
		\item Si $H \normsub K$ y $K \normsub G$ entonces $H \normsub G$
		\item Si $H \normsub G$ y $|H| = m$ entonces $H$ es el único subgrupo de $G$ de orden $m$.
		\item Si $H \normsub G$ entonces $H < Z(G)$.
		\begin{proof}[FALSO]
			Contraejemplo: En $G = D_4$ tomamos $H = \gen{B^2} = \{1, B, B^2, B^3\} \not\subset Z(D_4) = \{1, B^2\}$.
		\end{proof}
	\end{enumerate}
\end{ex}

\begin{ex}[H2.10]
	\begin{proof}
		Fijado $n$ y definida $\alpha_n : G \to G,\ x \mapsto x^n$ tenemos que $\alpha_n$ es un homomorfismo de grupos. Además podemos expresar $H_2 = \ker \alpha_n \implies H_2 \normsub G$. Además también tenemos que $H_1 = \ima \alpha_n < G$. Veamos que $H_1 \normsub G$. Es decir, que $gH_1\inv{g} = H_1,\ \forall g \in G$. Para ello tomamos $x_1^n \in H_1$ y lo conjugamos $gx_1^n\inv{g} = (gx_1\inv{g})^n$ por ser $\alpha$ homomorfismo de grupos. En particular $(gx_1\inv{g})^n \in \ima \alpha \implies (gx_1\inv{g})^n \in H_1 \implies (gx_1\inv{g})^n = x_2^n$ para algún $x_2 \in H_1 \implies H_1 \normsub G$.
	\end{proof}
\end{ex}

\begin{ex}[H2.13] Si $A$ es un grupo abeliano con $n$ elementos y $k$ es un entero primo con $n$, demostrad que la aplicación $\varphi : A \to A$ definida por $\varphi(a) = a^k$ es un isomorfismo.
	\begin{itemize}
		\item $\varphi$ homomorfismo de grupos.
		\begin{proof}
			\begin{align*}
				\varphi(a)\varphi(b) = a^kb^k = (ab)^k = \varphi(ab)
			\end{align*}
		\end{proof}
		\item $\varphi$ biyectiva $\iff \varphi$ inyectiva ya que dominio y codominio coinciden
		\begin{proof}
			$\ker \varphi = \{a \in A \mid \varphi(a) = a^k = 1\}$. Probaremos que $a^k = 1 \iff a = 1$ y por tanto que $\ker \varphi = \{1\} \implies \varphi$ inyectiva. Sabemos que $a^k = 1 \iff o(a^k) = 1$. Sea $t = o(a) \divides n$. Distinguimos dos casos
			\begin{itemize}
				\item Si $t = 1$ entonces $a = 1$ y ya está
				\item Si $t > 1$ entonces $o(a^k) = \frac{t}{mcd(k, t)} = \frac{t}{1} > 1$ contradicción. Luego necesariamente $t = o(a) = 1$.
			\end{itemize}
		\end{proof}
	\end{itemize}
\end{ex}

\begin{ex}[H2.22]
	\label{ex:h2.22}
	Demostrad que si $G$ es un grupo no conmutativo y tiene orden $p^3$ ($p$ un número primo) entonces $Z(G)$ tiene orden $p$.
	
	\begin{proof}
		Sabemos que $Z(G) < G \implies |Z(G)| \divides |G| \implies |Z(G)| \in \{1, p, p^2, p^3\}$
		\begin{itemize}
			\item $|Z(G)| \neq p^3$ porque en tal caso $G$ sería conmutativo
			\item $|Z(G)| \neq 1$ porque $G$ es un p-grupo y por tanto su centro no es el trivial.
			\item Si $|Z(G)| = p^2$ entonces $|G/Z(G)| = p \implies G/Z(G)$ es cíclico lo que no es posible si $G$ no es abeliano.
		\end{itemize}
	
		Por descarte concluimos que $|Z(G)| = p$.
	\end{proof}
\end{ex}


\begin{ex}[H2.25]
	Sabemos que $\autom{\Z/12\Z} \isom \uds{\Z^\star/12\Z}$ donde $\Z^\star/12\Z$ es el grupo multiplicativo $(\{\overline{1}, \overline{2}, \overline{3}, \dots,  \overline{11}\}, \cdot)$. Queda $\uds{\Z^\star/12\Z} = (\{\overline{1}, \overline{5}, \overline{7}, \overline{11}\}, \cdot )$ y además da la casualidad que $\forall x \in \uds{\Z^\star/12\Z},\ o(x) = 2$ (todos los elementos son su propio inverso) por lo que no tenemos restricciones al definir $f: \Z/2\Z \to \autom{\Z/12\Z}$:
	\begin{align*}
		f:\Z/2\Z & \to \autom{\Z/12\Z} \isom \uds{\Z^\star/12\Z} \\
		e = \overline{0} &\mapsto 1 \\
		\overline{1} &\mapsto \{\overline{1}, \overline{5}, \overline{7}, \overline{11}\}
	\end{align*}
\end{ex}

\begin{ex}[H2.26]
	Sea $|G_1| = m, |G_2| = n,\ mcd(m, n) = 1$. Si $f:G_1 \to G_2$ es h. de g. sabemos que $o(f(a)) \divides o(a),\ \forall a \in G_1$. Además $o(a) \divides m \land o(f(a)) \divides n$ por el teorema de Lagrange (\ref{thm:lagrange}).
	\begin{align*}
		\begin{cases}
		o(a) \divides m \land o(f(a)) \divides n \\
		mcd(m, n) = 1 \\
		o(f(a)) \divides o(a)
		\end{cases} \implies o(a) = o(f(a)) = 1, \forall a \in G_1
	\end{align*}
	Por lo que solo puede haber un homomorfismo entre ellos y además es el trivial $f(a) = e_{G_2}$.
\end{ex}

\begin{ex}[H2.19]
	Definimos una función $f:[0, 2\pi] \subset \R \to \mathbb{S}^1, \alpha\mapsto \cos \alpha + i \sin \alpha$. Esta función tiene la propiedad de que $f(\alpha) \cdot f(\alpha') = \cos(\alpha + \alpha') + i\sin(\alpha + \alpha') = f(\alpha + \alpha')$ y por tanto es un h. de g.\footnote{A la izquierda (en $\R$) sumamos pero a la derecha (en $\mathbb{S}^1$) multiplicamos.} entre $R$ y $\mathbb{S}^1$.
	
	Un elemento de $\cos \alpha + i \sin \alpha \in \mathbb{S}^1$ es de torsión $\iff \exists n \mid (\cos \alpha + i \sin \alpha )^n = 1$. Ahora bien $(\cos \alpha + i \sin \alpha )^n = \cos n\alpha + i \sin n\alpha = 1 \iff n\alpha = k 2\pi$. 
\end{ex}

\section{Hoja 4}

\begin{ex}[H4.11] Hallar los subgrupos de Sylow de $S_5$.
	Sabemos que $|S_5| = 5! = 2^3\cdot3\cdot 5$ y por \nameref{thm:sylow1} tenemos lo siguiente:
	\begin{itemize}
		\item $\exists P_2,\ |P_2| = 2^3 = 8$.
		\item $\exists P_3,\ |P_3| = 3$. Además en $S_5$ hay $\binom{5}{3}2! =20 $ 3-ciclos y en cada $gP_3\inv{g} = \{1, a, a^2 \mid o(a) = 3\}$ hay 2 elementos de orden 3 distintos. Además, $g_1P_3\inv{g_1} \cap g_2P_3\inv{g_2} = \{e\}$ porque si su intersección fuera más grande entonces serían el mismo subgrupo (porque son cíclicos). Es por esto que tenemos que repartir 40 elementos dando 2 a cada 3-grupo con lo que obtenemos $n_3 = 20 / 2 = 10$ 3-subgrupos de Sylow en $S_5$.
		\item $\exists P_5,\ |P_5| = 5$. Además en $S_5$ hay $4! = 24$ 5-ciclos (elementos de orden 5) y en cada $gP_5\inv{g} = \{1, a, a^2, a^3, a^4 \mid o(a) = 5\}$ tenemos $4$ elementos de orden $5$ distintos. Además, $g_1P_5\inv{g_1} \cap g_2P_5\inv{g_2} = \{e\}$ porque si su intersección fuera más grande entonces serían el mismo. Así, tenemos $24$ 5-ciclos a repartir entre los diferentes $gP_5\inv{g}$ dando 4 5-ciclos a cada 1. Por tanto tenemos $n_5 = 24/4 = 6$ 5-subgrupos de Sylow en $S_5$.
	\end{itemize}
\end{ex}

\begin{ex}[H4.18]
	Demostrar que todo grupo de orden $|G| = 5^3 \cdot 7^3$ tiene un subgrupo normal de orden 125.
	
	\begin{proof}
		\nameref{thm:sylow1} $\implies \exists P_5 < G,\ |P_5| = 125$. \nameref{thm:sylow3} $\implies n_5 \divides 7^3 \land n_5 \equiv 1 \mod 5$ es decir $n_5 \in \{1, 7, 49, 343\} \land n_5 \in \{1, 6, 11, \dots\}$. Como ni $49$ ni $343$ son conguentes con $1$ módulo 5 tenemos que $n_5 = 1 \implies P_5 \normsub G$.
	\end{proof}
\end{ex}

\begin{ex}[H4.20] Hallad todos los grupos abelianos de órdenes 36, 64, 96 y 100.
	\begin{enumerate}
		\item $|G| = 36 = 2^2\cdot3^2$
		\begin{proof}
			\nameref{thm:sylow1} $\implies \exists P_2,\ |P_2| = 4 \land \exists P_3,\ |P_3| = 9$. Además $G$ abeliano $\implies P_2, P_3 \normsub G \implies G \isom P_2 \times P_3$. Estudiamos los grupos de orden 4 y de orden 9
			\begin{itemize}
				\item $|P_2| = 4$ entonces $P_2 \isom \Z/2\Z \times \Z/2\Z \lor P_2 \isom \Z/4\Z$
				\item $|P_3| = 9$ entonces $P_3 \isom \Z/3\Z \times \Z/3\Z \lor P_3 \isom \Z/9\Z$
			\end{itemize}
			Como $G \isom P_2 \times P_3$ tenemos 4 posibles grupos abelianos de orden 36.
		\end{proof}
	\end{enumerate}
\end{ex}

\begin{ex}[H4.22]
	Hallar todos los grupos abelianos de orden 175.
	
	\begin{proof}
		$|G| = 5^2\cdot 7$. Por el \nameref{thm:sylow1} tenemos que $\exists P_5, P_7 < G$ con $|P_5| = 25,\ |P_7| = 7$ y además por ser $G$ abeliano tenemos que $P_5, P_7 \normsub G \implies G \isom P_5 \times P_7$. Estudiamos los grupos de orden 25 y de orden 7:
		\begin{itemize}
			\item $|P_5| = 25 \land P_5$ abeliano $\implies P_5 \isom \Z/25\Z \lor P_5 \isom \Z/5\Z \times \Z/5\Z$. En ambos casos $P_5$ es producto directo de cíclicos pues $\Z/n\Z$ es cíclico.
			\item $|P_7| = 7 \land P_7$ abeliano $\implies P_7 \isom \Z/7\Z$. Ocurre lo mismo que con $P_5$.
		\end{itemize}
		Concluimos que $G \isom \Z/25\Z \times \Z/7\Z \lor G \isom \Z/5\Z \times \Z/5\Z \times \Z/7\Z$. Los dos casos son abelianos por ser producto directo de grupos cíclicos.
	\end{proof}
\end{ex}

\begin{ex}[H4.23]
	¿Cuántos elementos de orden 3 puede tener un grupo abeliano de orden 36?
	
	\begin{proof}
		$|G| = 36 = 2^2 3^2$. \nameref{thm:sylow1} $\implies \exists P_2, P_3 < G,\ |P_2| = 4,\ |P_3| = 9$. Estudiamos los grupos de órdenes 4 y 9:
		\begin{itemize}
			\item $|P_3| = 9 \implies P_3 \isom \Z/9\Z \lor P_3 \isom \Z/3\Z \times \Z/3\Z$
			% TODO: terminar
		\end{itemize}
	\end{proof}
\end{ex}

\chapter{\'{I}ndices}

%%TODO por qué no podemos utilizar only= para hacer whitelisting en lugar de blacklisting? que alguien lo arregle parfavar
\renewcommand{\listtheoremname}{Lista de definiciones}
\listoftheorems[ignore={thm,eg,pro,cor,obs,lm,ex,th_ex}]

\renewcommand{\listtheoremname}{Lista de teoremas}
\listoftheorems[onlynamed,ignore={dfn,eg,cor,obs,ex,th_ex}]

\renewcommand{\listtheoremname}{Lista de ejemplos}
\listoftheorems[onlynamed,ignore={dfn,thm,pro,cor,obs,lm,ex,th_ex}]

\renewcommand{\listtheoremname}{Lista de ejercicios}
\listoftheorems[ignore={dfn,thm,pro,cor,obs,lm,eg}]

\printglossaries

\bibliographystyle{alpha}

\end{document}
